% ---------------------------------------------
% Copyright Ben Paul Wise. All rights reserved.
% ---------------------------------------------
% process with TeXworks as pdfLaTeX
%
% Survey of prior work bearing on acv-preliminary.tex, with a numerical
% comparison against the lasso and stochastic search variable selection.
% Numbers produced by examples/compete.py and examples/ssvs_sensitivity.py

\documentclass[english, 11pt]{report}

% ---------------------------------------------
% Copyright Ben Paul Wise. All rights reserved.
% ---------------------------------------------
% process with TeXworks as pdfLaTeX+MakeIndex+BibTex


% This file is \input by main.tex as the STYLE MODULE, so the class line is
% commented out here and the \documentclass (report, for chapters) lives in
% main.tex; everything below applies unchanged.
%\documentclass[english, 10pt]{article}
% `english' is necessary for babel, not `letter' or `a4paper'
% default is 10pt unless specified otherwise.
% I reserve 12 pt for legible final documents


%\usepackage[a4paper]{geometry}
\usepackage[letterpaper]{geometry}

\usepackage{geometry}
\geometry{verbose,lmargin=2.0cm,rmargin=2.0cm}

\setlength{\parindent}{2em}
\setlength{\parskip}{0.75em}

\usepackage{fancyhdr}
\pagestyle{fancy}
\usepackage{color}
\definecolor{mypurple}{RGB}{128, 0, 128}
\definecolor{myTeal}{RGB}{0, 128, 128}
\definecolor{myDarkGreen}{RGB}{0, 92, 0}
\definecolor{myLightGreen}{RGB}{0, 192, 0}
\definecolor{myDarkRed}{RGB}{192, 0, 0}
\definecolor{myDarkBlue}{RGB}{0, 0, 192}
\definecolor{myBeige}{RGB}{255, 255, 192}

\definecolor{GroupWBlue1}{RGB}{75, 99, 175} 
\definecolor{GroupWGrey1}{RGB}{145, 147, 150} 

\definecolor{schrift}{RGB}{0,73,114}

\usepackage{amsmath} %% get AMS standards for theorems, e.g. the hollow QED box,  \qedsymbol
%% But I like the solid "tombstone", Halmos, whatever:
\newcommand{\qedsolid}{\rule{1ex}{1.62ex}}
\usepackage{amssymb} % \sum


\usepackage{empheq} % colored equation-boxes
\usepackage[most]{tcolorbox}
\tcbset{colback=white!90!yellow, % i.e. 90% and 10% yellow
        colframe=green!36!blue, % i.e. 64% green and 36% blue, `schrift'
        highlight math style= {enhanced, %<-- needed for the ’remember’ options
            colframe=red,colback=red!10!white,boxsep=0pt}
        }
		
		
% setup colored and boxed quotes	
% this would give a yellowish background inside a reddish frame:
% colback=yellow!10!white, colframe=red!50!black, 

%\usepackage[most]{tcolorbox}     
\newtcolorbox{myquote}{
colback=yellow!10!white,
colframe=blue!50!black,
grow to right by=-20mm,
grow to left by=-30mm,
%boxrule=2pt,
boxsep=0pt,
breakable} \makeatletter

\newcommand{\blockquote}[1]{  \begin{myquote}  #1  \end{myquote}  }


% These are usually 'schrift' not 'GroupWBlue1'
\usepackage{hyperref}   % clickable links into or outfrom this text
\hypersetup{
    colorlinks=true,
    linkcolor=GroupWBlue1,
    filecolor=magenta,      
    urlcolor=cyan,
    citecolor=myTeal
}

\numberwithin{equation}{section} % number equations as (sec.eq), not just (eq)

% Good reference on tables in LaTeX:
% https://en.wikibooks.org/wiki/LaTeX/Tables 

\usepackage{graphicx}
% parameters for the 'figure box'
\fboxsep=1mm %padding thickness
\fboxrule=2pt %border thickness

\usepackage{sectsty}

% These are usually 'schrift' not 'GroupWBlue1' or 'GroupWGrey1'
\sectionfont{\color{GroupWBlue1}}  % sets colour of sections
\subsectionfont{\color{GroupWBlue1}}
\subsubsectionfont{\color{GroupWBlue1}}

\usepackage{lastpage}

% The following two are needed to shade individual cells,
% without the `feature' of over-writing the cell boundaries.
% However, they seem to clash with empheq, so they are commented out.
% this means that \cellcolor{color}{stuff} cannot be used; see revisions just prior to 2022-03-27 for colored cells.
%\usepackage[table]{xcolor}% http://ctan.org/pkg/xcolor, https://tex.stackexchange.com/questions/50349/color-only-a-cell-of-a-table
%\usepackage{hhline}% http://ctan.org/pkg/hhline, https://tex.stackexchange.com/questions/25062/cell-shading-without-removing-lines


\usepackage[T1]{fontenc} % necessary for babel
\usepackage{babel}
\usepackage{url} % so bitex can handle URL


%\rhead{Basic CP}   % default rhead is section
%\lhead{B. P. Wise} % default lhead is subsection

% Some people say that ragged right edges make it easier to read.
\raggedright

% ---------------------------------------------
% Copyright Ben Paul Wise. All rights reserved.
% ---------------------------------------------


% The report class numbers sections as chapter.section.  This document uses
% no chapters, so the chapter counter stays at zero and the sections come
% out as 0.1, 0.2 and so on.  Numbering them by section alone gives 1, 2, 3.
\renewcommand{\thesection}{\arabic{section}}

\def\myversion{v01a}

\lfoot{\color{myDarkRed}\bf{DRAFT \myversion}}
\cfoot{{\thepage} of {\pageref*{LastPage}}}
\rfoot{\copyright \  Ben P. Wise}

\begin{document}

\begin{titlepage}
\title{Prior Work Bearing on the Recovery of a Sparse Matrix\\
       A Survey of the Open Literature, with a Numerical Comparison}
\author{Ben Paul Wise}
\date{\today}

\maketitle

\begin{abstract}
This document reports a survey of the published literature bearing on the
companion paper \emph{Recovering a Sparse Matrix from Paired Observations}. Its
purpose is to establish which of the results in that paper are already known, so
that none of them is presented as new by mistake. The survey finds that the great
majority of the technical content is standard material carefully assembled, that
two of the findings restate results available in the standard references, and that
the recovery of a block arrangement from data is the subject of an active
literature. Four items were not found in print. The two methods a reader is most
likely to name as alternatives, the lasso and stochastic search variable selection,
are described at enough length to orient a reader who has not used them, and both
are run on the two illustrations of the paper. Every reference carries a link that
was checked and resolved.
\end{abstract}

\end{titlepage}

\clearpage

\raggedright

\section{Purpose and method of the survey}

The companion paper contains several results that appear interesting. Before any of
them is shown to colleagues, it is necessary to know which are already published,
including obscurely, so that no mistaken claim of novelty is made in front of an
audience.

The paper was decomposed into ten technical elements, and each was searched
separately rather than searching for the paper as a whole. The elements are the
model, the two bounded priors, the pattern prior, the hyper-prior on the occupancy
rate, the integrated criterion, the search procedure, the underdetermined regime,
the recovery of block arrangement, and the AutoClass framing. Searching for them
together returns nothing, because no single publication contains the same
combination. Searching for them separately returns a great deal.

Every reference given below was retrieved and resolved at the time of writing.
Where a publisher blocks automated access, an open copy or a stable repository
record is given instead.

\section{Verdict at a glance}

\begin{center}
\footnotesize
\fbox{
\begin{tabular}{lll}
Element in the paper & Status & Nearest prior work \\
\hline
The model, with the pattern unknown & well established & \cite{gsn2008}, \cite{basu2015} \\
Each row a separate subset problem & well established & \cite{silander2006} \\
Bounded log-uniform prior on $|A_{ij}|$ & known, different name & \cite{feroz2011} \\
Failure of the improper coefficient prior & textbook & \cite{llorente2022}, \cite{kass1995} \\
Failure of the improper prior on $\sigma$ & same result again & \cite{llorente2022} \\
Bernoulli pattern prior with Beta hyper-prior & well established & \cite{scott2010} \\
The increment growing only as $\log \beta$ & consistent with theory & \cite{scott2010} \\
Exponent $(n-M)/2$ and the Occam term & textbook & \cite{schwarz1978}, \cite{mackay1992} \\
Maximising within, integrating across & common practice & \cite{kass1995} \\
Fixed-$\sigma$ dynamic programming over rows & not found in this form & \cite{silander2006} \\
Block arrangement recovered from data & actively published & \cite{gudmundsson2021} and four others \\
AutoClass as the parent formalism & not found & \cite{cheeseman1996} \\
Attribution of over-fitting to the truncation & not found & --- \\
\end{tabular}
}
\end{center}

\section{The parts that are already well known}

These should be stated as received results. Presenting any of them as new will draw
a correction.

\subsection{The model and the estimation problem}

With $X$ known and $A$ unknown, the $i$th row of $A$ is an ordinary linear
regression of the $i$th row of $Y$ on the rows of $X$, with $T$ observations and
$N$ candidate coefficients. The whole problem is $N$ separate subset-selection
problems sharing one scale. This is among the most heavily worked problems in
modern statistics.

The Bayesian treatment by indicator variables goes back to George and McCulloch,
and was carried to vector autoregressions by George, Sun and Ni
\cite{gsn2008}, which is the closest published antecedent to the present
formulation. That paper places a two-component normal prior on each coefficient, a
Bernoulli prior on each indicator, and searches by Markov chain Monte Carlo. The
differences from the present work are the shape of the coefficient prior, the use
of a maximum rather than a sample, and the search. The structure of the argument is
the same. Section \ref{sec:ssvs} below describes the method at length, and
Section \ref{sec:results} reports what it does on the illustrations of the paper.

The regularised treatment of the same problem is due to Basu and Michailidis
\cite{basu2015}, with structured extensions in \cite{nicholson2015}. The same
problem is standard in signal processing, where $X$ is called a dictionary and the
several columns are called multiple measurement vectors; the Bayesian treatment
there is called sparse Bayesian learning \cite{zhang2011}, \cite{ament2021}.

\subsection{Why the improper priors fail}

Section 4.2 of the paper shows that $1 / |a|$ over $(0,\infty)$ leaves the
posterior without a maximum, and Section 5.1 shows that $1 / \sigma$ over
$(0,\infty)$ makes an exact fit infinitely good. Both are instances of one known
result: an improper prior on a parameter that appears in some candidate models and
not in others leaves the comparison undefined, because the missing normalising
constant does not cancel. This is the Bartlett--Lindley paradox, and it appears in
every treatment of Bayesian model selection \cite{llorente2022}, \cite{kass1995},
\cite{lambert2019}.

The derivations in the paper are correct and are worth keeping for the reader, but
they should be introduced as instances of a known difficulty rather than as
discoveries.

\subsection{The bounded log-uniform prior}

Truncating $1 / |a|$ to $[a_{\min}, a_{\max}]$ and dividing by
$\log (a_{\max} / a_{\min})$ is a known construction. In the astronomical
literature it is called the modified Jeffreys prior and is used for exactly the
reason given in Section 4.2 of the paper: the scale-free form is the right shape
but is not a density, and a model comparison needs a density \cite{feroz2011}. The
observation that $L_{A}$ is unchanged by a rescaling of units is the standard
motivation. The more common remedy in statistics is a mixture of $g$-priors
\cite{liang2008}.

One qualification. Although the device is standard, the criterion that follows from
it, carrying a term $\sum \log |A_{ij}|$ that grows as the coefficients grow, is
rarely written out, because most authors use a normal or a double-exponential
prior, whose corresponding term shrinks as the coefficients grow. The criterion was
not found displayed in the form of equation (7.1) anywhere. That is a point about
exposition and not about method, and it should be claimed no more strongly.

\subsection{The Beta hyper-prior and multiplicity}

This is the closest match in the entire survey, and it is a close one. Scott and
Berger \cite{scott2010} is the standard reference for placing a Beta prior on the
occupancy rate of a Bernoulli pattern prior and studying what that does to the
number of coefficients selected. The result in Section 6.2 of the paper, that the
cost of one more occupied cell is

\begin{equation}
\Delta \;=\; \log \frac{q - M - 1 + \beta}{M + \alpha}
\label{eq:increment}
\end{equation}

and that this grows only as $\log \beta$, is the log-odds increment of the
beta-binomial distribution. It is correct, it is worth displaying, and it is
consistent with the published theory rather than contrary to it.

The claim to avoid is that the hyper-prior was found to be unable to control
over-fitting. What Scott and Berger establish is the complementary and more precise
statement: the Beta prior supplies multiplicity correction, which is a different
thing from the Occam penalty that comes from integrating over the coefficients, and
neither substitutes for the other. The paper reaches the same conclusion by a
direct calculation. The right course is to cite them and say so.

\subsection{The integrated criterion}

Section 9 of the paper integrates the coefficients out by a Laplace approximation,
obtains the volume term, and finds the residual exponent falling from $n/2$ to
$(n-M)/2$. Both are textbook. The exponent $(n-M)/2$ is the ordinary residual
degrees of freedom of a normal linear model with the scale integrated out, and the
determinant term is the Occam factor \cite{schwarz1978}, \cite{mackay1992},
\cite{kass1995}. Retaining the constant terms rather than dropping them, as the
Schwarz criterion does, is the only departure, and it is a departure toward the
exact calculation rather than away from it.

The related account in which $\sum \log |A_{ij}|$ and the Occam term appear as code
lengths is the minimum description length literature \cite{wallace2005},
\cite{dhillon2011}. It is worth one sentence in the paper, because a listener from
that tradition will recognise the criterion.

\subsection{The search}

Enumerating subsets row by row, then combining rows, exploits the fact that the
score is a sum of per-row terms once the shared parameters are fixed. That property
is called decomposability, and it is the same property that makes exact structure
learning of Bayesian networks possible by dynamic programming \cite{silander2006}.
Deterministic Bayesian subset selection of other kinds appears in
\cite{hybrid2020} and \cite{kowal2022}.

The specific device in Section 10.2 of the paper, conditioning on a grid of
$\sigma$ values so that a criterion which is not separable becomes separable,
running the dynamic programme at each grid point, and scoring the resulting
candidates exactly, was not found published. It is a small idea and a natural one,
so absence from the search is weak evidence. It should be described, not claimed.

\section{The two methods a reader will name}

Two alternatives will come up whenever this problem is described, and a reader who
has not used them needs enough of each to judge the comparison in
Section \ref{sec:results}. Both are given here in the notation of the paper, so
that one row of the matrix is a regression of the $T$-vector $y$ on the $T$ by $N$
matrix $Z = X'$ with the $N$-vector $a$ to be found.

\subsection{The lasso}
\label{sec:lasso}

The lasso \cite{tibshirani1996} replaces least squares by least squares with a
penalty proportional to the sum of the absolute values of the coefficients

\begin{equation}
\hat{a}(\lambda) \;=\; \arg\min_{a} \; \tfrac{1}{2} \| y - Z a \|^{2}
\;+\; \lambda \sum_{j=1}^{N} |a_{j}|
\label{eq:lasso}
\end{equation}

Two facts account for its position. The first is that the absolute value has a
corner at the origin, and a penalty with a corner drives coefficients to exactly
zero rather than merely close to it. Selection and estimation are therefore done in
one step, and no search over patterns is required at all. The second is that the
problem is convex, so it is solved reliably and quickly at sizes far beyond
anything in the paper. Together these make it the default tool for the regime in
which the paper works, and the reason a colleague will ask about it.

Its relation to the present work is closer than it looks. The minimiser of
\eqref{eq:lasso} is the maximum-posterior estimate under a double-exponential prior
on each coefficient, with $\lambda$ playing the part of the rate
\cite{park2008}. It is therefore an estimate of exactly the kind the paper uses,
differing in the prior rather than in the principle. The differences are that the
double-exponential is peaked at zero rather than bounded away from it, and that
there is no separate indicator for occupancy: a cell is empty when the penalty
happens to drive it to zero, not because a discrete choice was made about it.

Three properties matter for the comparison. The first is that the penalty must be
supplied. Cross-validation is the usual answer, but cross-validation is aimed at
prediction rather than at recovering the pattern, and it is known to select too
many cells; with six observations in a row there is also very little to
cross-validate upon. The second is that the surviving coefficients are pulled
toward zero, so their values are biased, and refitting by least squares on the
selected pattern is normal practice. The third is that the lasso recovers the
correct pattern only when the columns of $Z$ satisfy a condition relating the
occupied columns to the empty ones \cite{zhao2006}, and that condition cannot be
checked without knowing the answer.

The adaptive lasso \cite{zou2006} addresses the third point by penalising each
coefficient in inverse proportion to a first-pass estimate of its size, so that
large coefficients are penalised little and small ones a great deal. It recovers
the correct pattern under weaker conditions, and it is the strongest variant that
can fairly be asked for here.

\subsection{Stochastic search variable selection}
\label{sec:ssvs}

Stochastic search variable selection, due to George and McCulloch and carried to
vector autoregressions in \cite{gsn2008}, attaches an indicator $\gamma_{ij}$ to
each cell and makes the prior on the coefficient depend upon it

\begin{eqnarray}
A_{ij} \mid \gamma_{ij} = 0 &\sim& N(0, \tau_{0}^{2})
\label{eq:ssvs_A}
\\
A_{ij} \mid \gamma_{ij} = 1 &\sim& N(0, \tau_{1}^{2})
\label{eq:ssvs_B}
\\
\gamma_{ij} &\sim& \mathrm{Bernoulli}(p)
\label{eq:ssvs_C}
\end{eqnarray}

with $\tau_{0}$ small and $\tau_{1}$ large. The essential device is that neither
component is exactly zero. Because every coefficient always has a value, every
conditional distribution in the model is a standard one, and the whole posterior
can be explored by drawing in turn the coefficients, the indicators, the occupancy
rate $p$ and the error variance. No move that changes the number of parameters is
ever needed, which is what makes the sampler simple enough to be practical.

The output is not a single pattern but a sample of them. The proportion of draws in
which a cell is occupied is its posterior inclusion probability, and the usual
summary is the median probability model, which takes every cell whose probability
is at least one half. That choice is not arbitrary; it is optimal for prediction
under stated conditions \cite{barbieri2004}.

The resemblance to the paper is very close. The pattern prior is the same Bernoulli
form, the hyper-prior on $p$ is the same Beta, one error scale serves the whole
matrix, and the hierarchy is arranged in the same order. Three things differ. The
coefficient prior is a pair of normal densities rather than a bounded log-uniform.
The pattern is sampled rather than maximised, so the answer is a set of
probabilities rather than one matrix. And no integral is performed in closed form;
everything is done by Monte Carlo.

The two widths $\tau_{0}$ and $\tau_{1}$ play exactly the part played by
$a_{\min}$ and $a_{\max}$ in the paper. They are not determined by the data, the
analyst must supply them, and the answer moves when they move.
Section \ref{sec:sens} measures by how much.

\subsection{What size of problem this method has been applied to}
\label{sec:ssvsscale}

This question was asked because the sampler failed to converge at $N = 200$, and it
matters a great deal whether that is a defect of the implementation used here or a
property of the method at that size. The published record is clear on one point: no
application of stochastic search variable selection at anything approaching forty
thousand coefficients was found. Every figure below was read out of the paper itself
and not out of an abstract.

\begin{center}
\footnotesize
\fbox{
\begin{tabular}{llll}
study & model & coefficients & draws \\
\hline
\cite{gsn2008}, application & 7 variables, 12 lags & $595$ & $20000$ after $10000$ \\
\cite{gsn2008}, simulations & 4 variables, 2 lags & about $36$ & $50000$ after $10000$ \\
\cite{korobilis2013} & 3 equations, 168 regressors each & $504$ & $150000$ after $50000$ \\
\cite{chan2020} & 20 variables, quarterly & --- & --- \\
\end{tabular}
}
\end{center}

Two things stand out. George, Sun and Ni describe $595$ coefficients as a large
number of parameters, and give that as the reason for the number of cycles they ran.
The problem of Section 12 of the companion paper carries forty thousand, which is
sixty-seven times as many. And the number of sweeps for each coefficient differs by
two or three orders of magnitude: about $34$ in \cite{gsn2008}, about $298$ in
\cite{korobilis2013}, against $0.2$ at eight thousand sweeps and $0.5$ at twenty
thousand in the runs reported here.

\subsection{Whether slow mixing is to be expected}

It is, and there is theory saying so. Yang, Wainwright and Jordan \cite{yang2016}
prove that posterior concentration, which is the statistical property that makes the
spike-and-slab posterior correct, does not imply rapid mixing of the algorithm used
to explore it. They obtain a mixing time linear in the number of covariates, up to a
logarithmic factor, only after introducing a truncated sparsity prior and analysing a
particular Metropolis--Hastings algorithm, and they distinguish explicitly between
rapid mixing, where the time grows polynomially, and slow mixing, where it grows
exponentially.

The plain Gibbs sampler of George and McCulloch is not the algorithm for which rapid
mixing has been proved, and a correct posterior is no guarantee that it can be
reached. There is also a body of work that exists because the plain sampler does not
scale: Biswas, Mackey and Meng \cite{biswas2022} reduce the cost of an iteration from
order $n^{2} p$; and in genomics, where the counts of predictors are far larger, the
field moved to variational methods because sampling was too slow for the size of
modern arrays.

\subsection{What this means for the comparison}

The implementation used here is the textbook one. Nothing found suggests it is
wrong, it reproduced the nine-variable results to the last bit after the change of
matrix inverse, and at nine variables it tied the method of the paper exactly. It is
nevertheless the plain sampler, and the published applications are sixty to eighty
times smaller while using a hundred to a thousand times more sweeps for each
coefficient.

The claim that stochastic search variable selection fails at two hundred variables is
therefore not defensible. The defensible claim is narrower: the plain Gibbs sampler,
run for as many sweeps as were affordable, has not converged at that size, which is
consistent both with published practice and with the theory. A blocked sampler, the
scalable variant of \cite{biswas2022}, or the truncated-prior scheme of
\cite{yang2016} might do very much better, and none has been tried.

\section{The comparison on the illustrations of the paper}
\label{sec:results}

Both methods were run on the two illustrations of the paper, over the same forty
seeds, with the same matrices, the same inputs and the same measurement errors. The
comparison is therefore paired at the level of the seed and at the level of the
case, giving eighty paired comparisons in all. The script is
\texttt{examples/compete.py}

Four variants of the lasso were run, because a single variant would not settle the
matter. The first chooses the penalty by leave-one-out cross-validation, which is
what an analyst without an estimate of the noise would do. The second chooses it by
the Schwarz criterion, given the true error scale, which no analyst would have. The
third is the adaptive lasso, reweighted from the second. The fourth chooses the
penalty, row by row, by counting the pattern errors against the truth and taking the
penalty that makes them fewest; this requires the answer and is therefore not a
method but a bound on what any choice of penalty could achieve. All four have their
coefficients refitted by least squares on the pattern they select, so that they are
charged for the pattern rather than for the shrinkage.

Stochastic search variable selection was run with $\tau_{0} = 0.02$ and
$\tau_{1} = 1.5$, twenty thousand draws of which the first five thousand were
discarded. Those widths bear the same relation to the true magnitudes, drawn
log-uniformly on $[0.40, 3.00]$, that the bounds $[0.20, 5.0]$ of the paper do, so
neither method is given better information than the other.

\subsection{What the forty realisations show}

\begin{center}
\footnotesize
\fbox{
\begin{tabular}{llccccc}
case & method & correct of 27 & false & missed & exact & median $\hat\sigma$ \\
\hline
block   & mple       & $26.20$ & $1.57$ & $0.80$ & $14$ of $40$ & $0.091$ \\
block   & lasso-cv   & $20.85$ & $18.43$ & $6.15$ & $0$ of $40$ & $0.614$ \\
block   & lasso-bic  & $23.70$ & $18.27$ & $3.30$ & $0$ of $40$ & $0.065$ \\
block   & alasso-bic & $23.32$ & $8.28$ & $3.67$ & $0$ of $40$ & $0.092$ \\
block   & lasso-best & $19.27$ & $1.68$ & $7.72$ & $0$ of $40$ & $0.746$ \\
block   & ssvs       & $26.32$ & $1.50$ & $0.68$ & $10$ of $40$ & $0.092$ \\
\hline
scatter & mple       & $26.40$ & $1.40$ & $0.60$ & $17$ of $40$ & $0.089$ \\
scatter & lasso-cv   & $21.73$ & $19.30$ & $5.28$ & $0$ of $40$ & $0.430$ \\
scatter & lasso-bic  & $23.98$ & $17.00$ & $3.02$ & $0$ of $40$ & $0.065$ \\
scatter & alasso-bic & $23.57$ & $7.85$ & $3.42$ & $2$ of $40$ & $0.096$ \\
scatter & lasso-best & $19.68$ & $2.35$ & $7.33$ & $0$ of $40$ & $0.696$ \\
scatter & ssvs       & $26.00$ & $1.12$ & $1.00$ & $17$ of $40$ & $0.094$ \\
\end{tabular}
}
\end{center}

The first line of each half reproduces the table of Section 10.5 of the paper
exactly, which confirms that the procedure compared against is the procedure the
paper describes.

Counting a pattern error as a cell wrongly occupied or wrongly empty, the eighty
paired comparisons fall out as follows

\begin{center}
\footnotesize
\fbox{
\begin{tabular}{lccc}
method & worse than mple & better & two-sided sign test \\
\hline
lasso-cv   & $80$ & $0$ & $p = 2 \times 10^{-24}$ \\
lasso-bic  & $80$ & $0$ & $p = 2 \times 10^{-24}$ \\
alasso-bic & $77$ & $2$ & $p = 1 \times 10^{-20}$ \\
lasso-best & $79$ & $0$ & $p = 3 \times 10^{-24}$ \\
ssvs       & $24$ & $24$ & $p = 1.00$ \\
\end{tabular}
}
\end{center}

\subsection{The lasso is not competitive here}

No variant of the lasso comes close, and the variant given the answer does not come
close either. The best of them, the adaptive lasso, occupies more than eight cells
that should be empty and misses nearly four that should be occupied, against a
little over two errors in total for the procedure of the paper. Averaged over the
eighty comparisons the total number of pattern errors is $11.61$ for the adaptive
lasso and $9.54$ for the oracle variant, against $2.19$ for the paper.

The reason is not obscure. With six observations and nine candidates in each row,
the condition under which the lasso recovers the correct pattern
\cite{zhao2006} is unlikely to hold, and a single continuous penalty must serve two
purposes at once: it must be large enough to empty six cells and small enough to
retain three. The oracle variant shows that no choice of penalty does both, since
even with the answer in hand it must trade nearly eight missed cells for its low
false count. Cross-validation, which is aimed at prediction, fails in the opposite
direction and occupies about nineteen cells too many.

This settles the objection that the paper compares itself only against the
minimum-norm pseudo-inverse. It does not, now, and the comparison is not close.

\subsection{Stochastic search variable selection is a dead heat}

This is the more important result, and it must be reported plainly. Against
stochastic search variable selection the paper's procedure wins twenty-four of the
eighty comparisons, loses twenty-four, and ties thirty-two. The sign test gives
$p = 1.00$, which is as even a result as the arithmetic permits. The average number
of pattern errors is $2.19$ for the paper and $2.15$ for the sampler. Exact
recoveries number thirty-one of eighty against twenty-seven, and a test of that
difference on the twelve discordant cases gives $p = 0.388$. The estimated error
scale is $0.0900$ against $0.0925$, with the truth at $0.1000$.

The honest summary is that the criterion developed in the paper attains the accuracy
of the established method at this size, and does not exceed it.

Two differences are real and are worth stating. The first is the worst case. The
sampler goes wrong less often, exceeding five pattern errors in six of the eighty
comparisons against eleven for the paper's procedure, but when it does go wrong it
goes further: its worst trial misplaces twenty-one cells against a worst of ten. The
second is that
the paper's procedure is deterministic: it evaluates a closed-form criterion on
candidates produced by a dynamic programme, and it returns the same answer every
time it is run, whereas the sampler returns a distribution and can differ between
chains on the same data.

\subsection{Both methods depend on widths the analyst supplies}
\label{sec:sens}

The most natural objection to the paper is that the bounds $a_{\min}$ and
$a_{\max}$ are not determined by the data. The same objection applies with at least
equal force to the alternative. The two widths of the two-component prior were
varied over four settings, with two chains at each, on seed $20260726$

\begin{center}
\footnotesize
\fbox{
\begin{tabular}{lccll}
$\tau_{0}$ & $\tau_{1}$ & ratio & block & scatter \\
\hline
$0.005$ & $3.00$ & $600$ & $M = 2$, 25 missed & $M = 7$ or $27$, chain-dependent \\
$0.020$ & $1.50$ & $75$ & 27 correct, 0 or 1 false & exact \\
$0.050$ & $1.00$ & $20$ & exact & exact \\
$0.100$ & $0.60$ & $6$ & $M = 81$, every cell & $M = 81$, every cell \\
\end{tabular}
}
\end{center}

Outside a fairly narrow band the method fails, and it fails in both directions. When
the wide component is made wider the model collapses to nearly empty, which is the
Bartlett--Lindley effect described in Section 3.2 acting inside the sampler itself:
spreading the prior over a larger range lowers the marginal likelihood of every
occupied cell. When the narrow component is made wider it ceases to represent an
empty cell at all, and every cell is occupied.

This does not excuse the paper. It does establish that the requirement is not
peculiar to the formulation of the paper but is a property of the problem, and that
the two bounds of the paper have an advantage of form: they are stated in the units
of the coefficients, as the smallest and largest magnitude worth calling non-zero,
which is a quantity an analyst can be asked about. A pair of prior standard
deviations whose effect is not monotone is harder to elicit.

\section{Block arrangement recovered from data}

This is the part of the work most exposed to correction, and it deserves care.

The premise is that the global vector autoregression of the economics literature
fixes a block arrangement in advance, so that recovering one from the data would be
of interest. The premise is sound, and it is also the premise of a live and growing
literature. Several groups have published methods that infer group or block
membership from the data instead of assuming it, and they do so with consistency
proofs and at cross-sections far larger than nine
\cite{gudmundsson2021}, \cite{brownlees2017}, \cite{billio2019},
\cite{biclustered2012}, \cite{kim2026}.

Plainly stated: the ability to recover a block arrangement from data, without
imposing it, is neither new nor rare. A claim of novelty on that ground will be
corrected, probably by whoever in the audience follows the econometrics literature.

There is, however, a real and defensible distinction, and it is worth making
precisely. Every method just cited contains machinery for grouping, whether a block
model, a clustering prior or a spectral step. Each is built to find groups. The
method in the paper contains no such machinery at all. It treats every cell alike
and fits each row separately, and the block arrangement, when there is one, emerges
as a by-product of selecting cells one at a time.

The paper already says the honest thing about this, in Section 10.6 and in the
forty-trial comparison: block and scattered arrangements are recovered equally
well, the sign test gives $p = 0.345$, and no difference should be expected, since
the prior treats all cells alike. That result is the right thing to present. It is
evidence that no block assumption is needed for this kind of recovery, which is a
statement about the block assumption rather than a claim of a new capability.
Framed that way it is defensible. Framed as the recovery of block structure it is
not.

One suggestion. The methods cited are aimed at a different regime, with many series
and long samples, where the object is to summarise a large system by a few groups.
The method of the paper is aimed at the regime in which the sample is shorter than
the cross-section and the ordinary estimate does not exist at all. Saying which
regime is addressed will forestall most of the objection.

\section{The AutoClass framing}

Cheeseman and Stutz \cite{cheeseman1996} applied the same hierarchical
construction, a discrete structural choice, parameters within the structure, and an
approximated marginal likelihood to choose among structures, to the problem of
classification. The paper carries that construction to the selection of a sparsity
pattern in a linear model. No publication was found that makes this extension or
that cites AutoClass in a sparse-regression setting.

The lineage is a matter of presentation rather than of method, and the underlying
approximation is the ordinary Laplace one, so the framing is legitimate but should
be offered as a way of organising the argument rather than as a technical
contribution.

\section{What appears to be new}

Four items survived the search. None is a new method. Each is small, checkable, and
safe to state.

The first is the explicit criterion. Equations (7.1) and (9.4) of the paper, with
the constants retained and with the term $\sum \log |A_{ij}|$ that follows from a
bounded log-uniform prior, are not displayed in this form in anything found. The
ingredients are all standard; the assembled and normalised criterion is not in
print.

The second is the diagnosis of the residual over-fitting. Section 10.6 attributes
the surviving falsely occupied cells to the Laplace step, on the evidence that they
cluster at small magnitudes, which is where the approximation ignores the
truncation at $a_{\min}$ and treats the product of reciprocals as constant. No
comparable attribution was found in the literature. It is a modest, verifiable and
genuinely useful observation, and it is the strongest thing in the paper to present
as new.

The third is the deterministic treatment throughout. Nearly all of the Bayesian
work cited above samples. This work maximises and integrates in closed form, with a
deterministic search. Deterministic Bayesian subset selection exists, but not, so
far as the search shows, with this criterion or in this regime.

The fourth is the fixed-$\sigma$ dynamic programme, as discussed above. It is weak
evidence, but it was not found.

To these should be added something that is not a result but is worth stating. The
paper is self-contained. Every constant is normalised, every integral is verified
in Maxima, and every number a reader needs in order to reproduce the illustrations
is printed in the appendix. Very little of the cited literature can be reproduced
from the paper alone. That is a real service to a reader and it is safe to say so.

\section{What a colleague is most likely to raise}

Anticipating these will be worth more than any of the above.

The first objection was that the paper compares itself only against the minimum-norm
pseudo-inverse, which recovers all eighty-one cells and is not a serious competitor.
That objection is now answered by Section \ref{sec:results}, and answered
comfortably. It should be answered in the paper as well, by a table, and not left to
be raised from the floor.

The second is that the criterion attains but does not exceed the accuracy of
stochastic search variable selection. Somebody will ask what the paper buys. The
answer is not that it is more accurate, because it is not. The answer is that it is
deterministic, that its worst case over eighty trials was less than half as bad, and
that it produces a criterion that can be examined term by term, so that the source
of a failure can be identified. Section 10.6 of the paper does exactly that, and no
sampler yields such an account. Saying this plainly is much stronger than claiming
an advantage the numbers do not support.

The third is that the bounds must be supplied by the analyst, and that Section 10.6
shows the answer moving when $a_{\min}$ moves from $0.20$ to $0.05$. This is a
genuine sensitivity, and the paper is candid about it. It is now known that the
alternative is at least as sensitive to its own two widths, and fails in both
directions outside a narrow band, which makes the objection one about the problem
rather than about the paper.

The fourth is the scale of the illustration. Nine by nine, with six observations,
against a literature that works at hundreds of series. Saying at the outset that
this is the smallest problem exhibiting the phenomenon, and that the size is chosen
so a reader can check every number, will keep the point from being pressed. It
should be added that the search of Section 10.2 enumerates subsets, and that nothing
in the paper establishes how the method behaves when $q$ is forty thousand.

The fifth is that the Laplace step is an approximation, and that Section 10.6 shows
it is the one that costs something. This is already acknowledged.

\section{Summary judgement}

Nothing in the paper duplicates a specific published result outright, and nothing
in it is a mistaken statement of what is known. But the great majority of its
technical content is standard material assembled carefully, and two of its
findings, the failure of the improper priors and the behaviour of the Beta
hyper-prior, restate results that are in the standard references.

Presented as an exposition of how a maximum-posterior criterion for a sparsity
pattern is built from first principles, with everything normalised and verified,
the work is sound and useful and will draw no correction. Presented as a new method
for recovering sparse or block-structured matrices, it will.

The numerical comparison sharpens this. Against the lasso, in four variants, two of
which are given information no analyst would have, the method of the paper wins
seventy-nine or eighty of eighty paired trials. Against stochastic search variable
selection it wins twenty-four, loses twenty-four and ties thirty-two. The first
result is worth presenting and closes the most likely objection. The second is worth
presenting because it is true, and because presenting it removes any suggestion that
the comparison was arranged to flatter.

Three things are worth putting forward. The diagnosis of where the residual
over-fitting comes from, which was not found in the literature. The demonstration
that no grouping mechanism is required in order to recover a grouped arrangement.
And the observation that a deterministic criterion, evaluated in closed form,
matches a Markov chain method at this size with a smaller worst case. All three are
modest, all three are supported by the numbers in hand, and all three will hold up.

\begin{thebibliography}{99}

\bibitem{ament2021}
S. Ament and C. Gomes,
\emph{Sparse Bayesian learning via stepwise regression},
Proceedings of Machine Learning Research 139, 2021.
\url{https://proceedings.mlr.press/v139/ament21a/ament21a.pdf}

\bibitem{biswas2022}
N. Biswas, L. Mackey and X.-L. Meng,
\emph{Scalable spike-and-slab}, 2022.
\url{https://arxiv.org/abs/2204.01668}

\bibitem{chan2020}
J. C. C. Chan,
\emph{Large Bayesian vector autoregressions}.
\url{https://www.joshuachan.org/papers/large_BVAR.pdf}

\bibitem{korobilis2013}
D. Korobilis,
\emph{VAR forecasting using Bayesian variable selection},
Journal of Applied Econometrics 28(2), 204--230, 2013.
\url{https://mpra.ub.uni-muenchen.de/21122/1/ssvs.pdf}

\bibitem{yang2016}
Y. Yang, M. J. Wainwright and M. I. Jordan,
\emph{On the computational complexity of high-dimensional Bayesian variable
selection},
Annals of Statistics 44(6), 2497--2532, 2016.
\url{https://projecteuclid.org/journals/annals-of-statistics/volume-44/issue-6/On-the-computational-complexity-of-high-dimensional-Bayesian-variable-selection/10.1214/15-AOS1417.pdf}

\bibitem{barbieri2004}
M. M. Barbieri and J. O. Berger,
\emph{Optimal predictive model selection},
Annals of Statistics 32(3), 870--897, 2004.
\url{https://projecteuclid.org/journals/annals-of-statistics/volume-32/issue-3/Optimal-predictive-model-selection/10.1214/009053604000000238.full}

\bibitem{basu2015}
S. Basu and G. Michailidis,
\emph{Regularized estimation in sparse high-dimensional time series models},
Annals of Statistics 43(4), 1535--1567, 2015.
\url{https://projecteuclid.org/euclid.aos/1434546214}

\bibitem{billio2019}
M. Billio, R. Casarin and L. Rossini,
\emph{Bayesian nonparametric sparse VAR models},
Journal of Econometrics 212(1), 97--115, 2019.
\url{https://arxiv.org/abs/1608.02740}

\bibitem{biclustered2012}
T.-Y. Chiang and others,
\emph{Learning bi-clustered vector autoregressive models},
ECML PKDD 2012, Lecture Notes in Computer Science 7524.
\url{https://link.springer.com/chapter/10.1007/978-3-642-33486-3_47}

\bibitem{brownlees2017}
C. T. Brownlees, G. S. Gudmundsson and G. Lugosi,
\emph{Sparse estimation of huge networks with a block-wise structure},
Econometrics Journal 20(3), S61--S85, 2017.
\url{https://ideas.repec.org/a/wly/emjrnl/v20y2017i3ps61-s85.html}

\bibitem{cheeseman1996}
P. Cheeseman and J. Stutz,
\emph{Bayesian classification (AutoClass): theory and results},
in Advances in Knowledge Discovery and Data Mining, 1996.
\url{https://ntrs.nasa.gov/citations/19920075544}

\bibitem{dhillon2011}
P. S. Dhillon, D. Foster and L. Ungar,
\emph{Minimum description length penalization for group and multi-task sparse
learning},
Journal of Machine Learning Research 12, 2011.
\url{https://jmlr.csail.mit.edu/papers/v12/dhillon11a.html}

\bibitem{feroz2011}
F. Feroz, S. T. Balan and M. P. Hobson,
\emph{Detecting extrasolar planets from stellar radial velocities using Bayesian
evidence}, 2011. Sets out the modified Jeffreys prior and the reason for it.
\url{https://arxiv.org/abs/1012.5129}

\bibitem{gsn2008}
E. I. George, D. Sun and S. Ni,
\emph{Bayesian stochastic search for VAR model restrictions},
Journal of Econometrics 142(1), 553--580, 2008.
\url{https://faculty.wharton.upenn.edu/wp-content/uploads/2012/04/GeorgeSunNi-JE2008.pdf}

\bibitem{gudmundsson2021}
G. S. Gudmundsson and C. T. Brownlees,
\emph{Detecting groups in large vector autoregressions},
Journal of Econometrics 225(1), 2--26, 2021.
\url{https://bse.eu/research/publications/detecting-groups-large-vector-autoregressions}

\bibitem{hybrid2020}
\emph{Bayesian selection of best subsets via hybrid search},
Computational Statistics, 2020.
\url{https://link.springer.com/article/10.1007/s00180-020-00996-y}

\bibitem{kass1995}
R. E. Kass and A. E. Raftery,
\emph{Bayes factors},
Journal of the American Statistical Association 90(430), 773--795, 1995.
\url{https://www.stat.washington.edu/raftery/Research/PDF/kass1995.pdf}

\bibitem{kim2026}
Y. Kim and C. Baek,
\emph{Latent community paths in VAR-type models via dynamic directed spectral
co-clustering}, 2026.
\url{https://arxiv.org/abs/2604.12563}

\bibitem{kowal2022}
D. R. Kowal,
\emph{Bayesian subset selection and variable importance for interpretable
prediction and classification},
Journal of Machine Learning Research 23, 2022.
\url{https://www.jmlr.org/papers/volume23/21-0403/21-0403.pdf}

\bibitem{lambert2019}
\emph{A parsimonious tour of Bayesian model uncertainty}. Surveys the
Bartlett--Lindley paradox and the standard remedies.
\url{https://arxiv.org/abs/1902.05539}

\bibitem{liang2008}
F. Liang, R. Paulo, G. Molina, M. A. Clyde and J. O. Berger,
\emph{Mixtures of g-priors for Bayesian variable selection},
Journal of the American Statistical Association 103(481), 410--423, 2008.
\url{https://www2.stat.duke.edu/~berger/papers/g-priors.pdf}

\bibitem{llorente2022}
F. Llorente, L. Martino, D. Delgado and J. L\'{o}pez-Santiago,
\emph{On the safe use of prior densities for Bayesian model selection}, 2022.
\url{https://arxiv.org/abs/2206.05210}

\bibitem{mackay1992}
D. J. C. MacKay,
\emph{Bayesian interpolation},
Neural Computation 4(3), 415--447, 1992.
\url{https://authors.library.caltech.edu/records/r7qgh-q6g10}

\bibitem{nicholson2015}
W. B. Nicholson, J. Bien and D. S. Matteson,
\emph{VARX-L: structured regularization for large vector autoregressions with
exogenous variables}, 2015.
\url{https://arxiv.org/abs/1508.07497}

\bibitem{park2008}
T. Park and G. Casella,
\emph{The Bayesian lasso},
Journal of the American Statistical Association 103(482), 681--686, 2008.
\url{https://people.eecs.berkeley.edu/~jordan/courses/260-spring09/other-readings/park-casella.pdf}

\bibitem{schwarz1978}
G. Schwarz,
\emph{Estimating the dimension of a model},
Annals of Statistics 6(2), 461--464, 1978.
\url{https://projecteuclid.org/journals/annals-of-statistics/volume-6/issue-2/Estimating-the-Dimension-of-a-Model/10.1214/aos/1176344136.full}

\bibitem{scott2010}
J. G. Scott and J. O. Berger,
\emph{Bayes and empirical-Bayes multiplicity adjustment in the variable-selection
problem},
Annals of Statistics 38(5), 2587--2619, 2010.
\url{https://arxiv.org/abs/1011.2333}

\bibitem{silander2006}
T. Silander and P. Myllym\"{a}ki,
\emph{A simple approach for finding the globally optimal Bayesian network
structure},
Proceedings of the Twenty-Second Conference on Uncertainty in Artificial
Intelligence, 2006.
\url{https://arxiv.org/abs/1206.6875}

\bibitem{tibshirani1996}
R. Tibshirani,
\emph{Regression shrinkage and selection via the lasso},
Journal of the Royal Statistical Society Series B 58(1), 267--288, 1996.
\url{https://webdoc.agsci.colostate.edu/koontz/arec-econ535/papers/Tibshirani%20%28JRSS-B%201996%29.pdf}

\bibitem{wallace2005}
C. S. Wallace,
\emph{Statistical and Inductive Inference by Minimum Message Length},
Springer, 2005.
\url{https://link.springer.com/book/10.1007/0-387-27656-4}

\bibitem{zhang2011}
Z. Zhang and B. D. Rao,
\emph{Sparse signal recovery with temporally correlated source vectors using
sparse Bayesian learning}, 2011.
\url{https://arxiv.org/abs/1102.3949}

\bibitem{zhao2006}
P. Zhao and B. Yu,
\emph{On model selection consistency of lasso},
Journal of Machine Learning Research 7, 2541--2563, 2006.
\url{https://www.jmlr.org/papers/volume7/zhao06a/zhao06a.pdf}

\bibitem{zou2006}
H. Zou,
\emph{The adaptive lasso and its oracle properties},
Journal of the American Statistical Association 101(476), 1418--1429, 2006.
\url{https://pages.stat.wisc.edu/~shao/stat992/zou2006.pdf}

\end{thebibliography}

\end{document}
